\algproblem{ALG01}{允许删除一次的最大子段和}
给定非空实数数组 $a_0,\ldots,a_{T-1}$。选取一个连续区间，可以从该区间删除至多一个元素，但至少保留一个。求保留元素之和的最大值，要求 $O(T)$ 时间、$O(1)$ 额外空间。
\solution 令 $f_i$ 为以 $i$ 结尾且不删除元素的最好值，$g_i$ 为以 $i$ 结尾、至多删除一个元素且仍非空的最好值；在 $g_i$ 中，位置 $i$ 本身可以被删。初始化 $f_0=g_0=a_0$。对 $i\ge1$，
\[
f_i=\max(a_i,f_{i-1}+a_i),\qquad
g_i=\max(a_i,g_{i-1}+a_i,f_{i-1}).
\]
最后一式依次对应重新开始、保留 $a_i$ 并延长已有区间、删除当前 $a_i$。若某个解删除了区间第一个元素，把区间起点后移即可得到相同保留序列，所以不需单独增加这种情况。其余合法解按末元素是否删除恰被递推覆盖。

答案为 $\max_i g_i$。更新时保存旧的 $f$，避免第三项误读新状态。始终以实际元素初始化，不能用零替代，否则全负数组会错误地返回空段。两状态的归纳证明同时保证最优性；时间 $O(T)$、额外空间 $O(1)$。

\algproblem{ALG02}{到达终点的最少跳跃次数}
非空数组 $a_0,\ldots,a_{T-1}$ 为非负整数。从位置0出发，在位置 $i$ 可向右跳 $1,\ldots,a_i$ 格。求到达 $T-1$ 的最少跳数；不能到达则返回无解。要求线性时间、常数额外空间，不能假设输入保证可达。
\solution $T=1$ 时返回0。维护当前层可达的最右位置 $R$、扫描到目前为止下一层可达的最右位置 $F$，以及已使用的层数 $k$，初始均为0。从 $i=0$ 扫到 $T-2$，先置
\[
F\leftarrow\max\{F,\min(T-1,i+a_i)\}.
\]
每当 $i=R$，当前层已全部扫描。若 $F=R<T-1$，则没有新的可达位置，返回无解；否则令 $k\leftarrow k+1$、$R\leftarrow F$。一旦 $R=T-1$ 就返回 $k$。

从一个位置允许跳到一整段前方位置，所以至多 $k$ 次跳跃能到达的位置构成从0开始的连续前缀。扫描完此前缀后，所有允许的一步延伸的并集恰到 $F$，因此新的边界是至多 $k+1$ 步的最远位置。算法逐层推进，首次覆盖终点的层数就是最短跳数。每个下标只扫描一次，时间 $O(T)$，额外空间 $O(1)$。

\algproblem{ALG03}{持有期限受限的一次交易}
价格 $p_0,\ldots,p_{T-1}$ 非负，给定整数 $1\le h\le L$。最多做一次交易，买入日 $i$ 与卖出日 $j$ 必须满足 $h\le j-i\le L$，每次卖出收取费用 $c\ge0$。可选择不交易。求最大净利润和一个最优日期对，要求 $O(T)$ 时间。
\solution 对卖出日 $j$，合法买入日构成
\[
\max(0,j-L)\le i\le j-h.
\]
若该区间为空，则当天没有可行交易。维护价格递增的单调队列：在第 $j$ 天把 $j-h$ 加入（若非负），入队时删去队尾价格严格更高的下标；再删除小于 $j-L$ 的队首。保留相等价格可选择较早买入日。

队首给出最低合法买价，候选利润为 $p_j-p_i-c$。答案从零开始，仅在严格改善时记录日期；若最终仍为零，返回不交易。若 $T\le h$，自然没有合法日期对。

固定卖出日时，费用与卖价固定，最低买价必定最优。队尾删除由更晚且更低的价格支配旧价格来保证；旧价格不会比新价格更晚过期。延迟加入保证最短持有期，队首过期保证最长持有期。每个下标只入队、出队一次，时间 $O(T)$，空间 $O(\min(T,L))$。

\algproblem{ALG04}{有限库存的逐日交易}
给定 $T$ 天的非负价格及整数库存上限 $K\ge0$。每天至多买入一单位、卖出一单位或不操作，禁止做空，买卖每单位均收手续费 $c\ge0$。初始现金和库存均为0，允许现金为负，最后库存必须为0。求最大终值，并给出 $O(T(K+1))$ 的动态规划。
\solution 令 $V_t(q)$ 为处理前 $t$ 天、持有 $q$ 单位时的最大现金，$0\le q\le K$。初值 $V_0(0)=0$，其余为 $-\infty$。在价格 $p_t$ 的一天，
\[
V_{t+1}(q)=\max\left\{
\begin{array}{l}
V_t(q),\\
V_t(q-1)-p_t-c\quad(q\ge1),\\
V_t(q+1)+p_t-c\quad(q<K).
\end{array}\right\}
\]
三项分别是不操作、买入、卖出。所有右端必须读取上一天的数组；若原地连续更新，可能错误地允许一天买卖多单位。

任意合法策略按最后一天的动作恰落入以上三类。去掉最后动作后得到相应旧状态，而每个有限旧状态接上该动作也合法，归纳证明递推同时给出上界与可达值。最终答案为 $V_T(0)\ge0$，因为始终不操作可行。时间 $O(T(K+1))$，滚动数组空间 $O(K+1)$；保存动作还可恢复交易日程。$K=0$ 或 $T=0$ 时答案都是0。允许现金为负是必要的模型说明，否则还需检查每笔买入是否有足够现金。

\algproblem{ALG05}{至多穿过一处障碍的路径计数}
一个 $R\times C$ 网格的每格标为0或1，1表示障碍。只能从左上角向右或向下移动到右下角，最多允许经过一个障碍格，起终点也按同一规则计入。求不同路径数对给定整数 $M\ge2$ 的余数，要求 $O(RC)$ 时间及 $O(C)$ 空间。
\solution 令 $A_{i,j}$、$B_{i,j}$ 分别计到达 $(i,j)$ 时经过零个、恰好一个障碍的路径数。先把上方和左方的两种计数分别相加，记为 $a,b$。若当前格为0，置 $(A,B)=(a,b)$；若为1，置 $(A,B)=(0,a)$。每次加法均取模 $M$。

起点单独初始化：起点为0则 $(A,B)=(1,0)$，为1则为 $(0,1)$。其余越界前驱一律贡献零。终点答案为 $A_{R-1,C-1}+B_{R-1,C-1}\pmod M$。

每条路径的倒数第二格恰为上方或左方，两类互斥；障碍计数又恰由当前格是否为1决定，所以递推既不遗漏也不重复。按行从左到右更新两个长度 $C$ 的数组，更新前的本列保存上方状态、已更新的前一列保存左方状态，先取出两者再覆盖即可。$R,C\ge1$ 时上述初始化也覆盖单格网格。

\algproblem{ALG06}{最短公共超序列}
给定字符串 $a,b$，允许空串。构造一个最短字符串 $s$，使 $a,b$ 都是 $s$ 的子序列。相同字符可以同时对应两串中的一个字符。求最短长度并恢复一个答案，给出正确性证明。
\solution 设前缀长度为 $i,j$ 时最短超序列长度为 $D_{i,j}$。边界是 $D_{i,0}=i$、$D_{0,j}=j$。若末字符相同，可让它们共用答案的最后一位：
\[
D_{i,j}=D_{i-1,j-1}+1\quad(a_i=b_j).
\]
这里字符从1编号。若末字符不同，则最后一位只能消去一个串的末字符，因此
\[
D_{i,j}=1+\min(D_{i-1,j},D_{i,j-1}).
\]
对相同末字符情形，任意超序列中两次末字符匹配都可改用较晚的那一位置，故共用最后一位不会变差；不同末字符的两个选择则穷尽了最优解。由前缀长度归纳，递推正确。

从 $(|a|,|b|)$ 回溯：相同则记录该字符并对角移动；不同则沿最小前驱记录被消去串的末字符；到边界后补入剩余前缀，最后反转记录。时间、存储均为 $O((|a|+1)(|b|+1))$。

若 $L$ 为最长公共子序列长度，则最短长度也等于 $|a|+|b|-L$：对齐两串时共用的位置形成公共子序列，最多共用 $L$ 位；反过来沿一条最长公共子序列，把两串间隙依次插入即可达到这个下界。

\algproblem{ALG07}{交织来源的计数与随机恢复}
给定字符串 $a,b,t$，其中 $|t|=|a|+|b|$。保持 $a,b$ 各自的顺序，把每个 $t$ 中的位置标为来自哪一个串。不同来源标记视为不同方案，即使字符相同。求合法方案数；若方案数非零，设计使所有合法标记等概率的随机恢复方法。
\solution 令 $C_{i,j}$ 为已经使用 $a$ 的前 $i$ 个字符、$b$ 的前 $j$ 个字符后，完成剩余 $t$ 的方案数。边界 $C_{m,n}=1$，其中 $m=|a|,n=|b|$。从后往前计算
\[
C_{i,j}=\ind_{\{i<m,\ a_{i+1}=t_{i+j+1}\}}C_{i+1,j}
       +\ind_{\{j<n,\ b_{j+1}=t_{i+j+1}\}}C_{i,j+1}.
\]
字符从1编号，越界分支贡献零。两项对应下一位置的不同来源，因此即使字符相同也应相加。总数是 $C_{0,0}$；若长度不符则直接为零。

恢复时处于 $C_{i,j}>0$ 的状态，按上述两个非负分支计数占总数的比例选择来源，再进入对应后继状态。任意完整标记的概率沿路径相乘，分子分母逐项相消，最终等于 $1/C_{0,0}$，故采样均匀。空串的唯一标记也被边界覆盖。

表含 $O((m+1)(n+1))$ 个状态，恢复恰走 $m+n$ 步。计数用精确整数；这里的状态运算次数界不把大整数加法假设为固定比特成本。只用模数计数将失去用于精确采样的权重。

\algproblem{ALG08}{数量受限的加权区间选择}
有 $n$ 个任务，任务 $i$ 占用半开区间 $[s_i,e_i)$，其中 $s_i<e_i$，收益 $w_i$ 可以为负。最多选择整数 $K\ge0$ 个互不重叠任务，允许一个任务结束时下一个立即开始。求最大总收益，要求给出 $O(n\log n+nK)$ 的算法。
\solution 按结束时间排序，记 $p(i)$ 为满足 $e_j\le s_i$ 的最大前驱下标，不存在则为0。由于 $s_i<e_i$，前驱一定排在 $i$ 之前；对有序结束时间二分即可求全部 $p(i)$，耗时 $O(n\log n)$。

令 $F(i,k)$ 是前 $i$ 个任务中至多选 $k$ 个的最好收益。边界 $F(0,k)=F(i,0)=0$，其余满足
\[
F(i,k)=\max\{F(i-1,k),\ w_i+F(p(i),k-1)\}.
\]
不选 $i$ 得到第一项；选 $i$ 时，其余所选任务必须在它开始前结束，故得到第二项。反过来第二项中的前驱最优集合与 $i$ 合并必定可行。这证明递推既是上界又可达到。

将 $K$ 截为 $\min(K,n)$，答案为 $F(n,K)$。负收益任务可被跳过，所以空集合的收益0始终是下界。完整表占 $O(n(K+1))$ 空间并支持回溯；按 $k$ 逐层计算，只求收益时可用两条长度 $n+1$ 的数组。

\algproblem{ALG09}{任务自带冷却期的日程选择}
第 $i$ 天有一项只能当天执行的任务，收益为实数 $r_i$。执行后必须休息接下来的 $L_i\ge0$ 个完整日，其中 $L_i$ 为整数；休息期间不能执行其他任务。任务可以跳过，$i=0,\ldots,T-1$。求最大总收益并恢复执行日期，要求 $O(T)$ 时间。
\solution 记 $V_i$ 为从第 $i$ 天起的最好收益，$V_T=0$。执行当天任务后，下一可执行日为
\[
q_i=\min(T,i+L_i+1).
\]
反向递推
\[
V_i=\max\{V_{i+1},r_i+V_{q_i}\},\qquad i=T-1,\ldots,0.
\]
第一项跳过当天，第二项执行并直接转入休息结束后的子问题。两种选择覆盖所有合法日程，而且拼接相应子问题的最优日程必合法，故由反向归纳得到最优性。

记录每个状态选取的分支。从0出发，跳过时前进一日，执行时输出 $i$ 并跳到 $q_i$；下标严格递增，所以恢复至多访问 $T$ 个状态。时间和空间均为 $O(T)$。若收益非正，可通过跳过得到不小于它的收益；$T=0$ 时答案0。由于 $q_i$ 的跨度各不相同，不能只保留固定数目的最近状态。

\algproblem{ALG10}{最少插入括号的修复}
给定只含左右圆括号的字符串，不能删除或改变原有字符，只能插入括号。求使整个字符串成为合法括号序列的最少插入数，并输出一种达到最少次数的修复。要求线性时间。
\solution 从左到右扫描，维护尚未配对的左括号数 $b$。遇到左括号，原样输出并令 $b$ 加1。遇到右括号：若 $b>0$，原样输出并令 $b$ 减1；若 $b=0$，先插入一个左括号，再输出该右括号，插入数加1。扫描结束后再补入 $b$ 个右括号。

输出的每个前缀左括号数不少于右括号数，末尾总数又相等，所以确为合法序列；原串的所有字符保持了顺序。

最优性可由前缀亏空证明。将空前缀也计入，记原串前缀余额的最小值为 $m\le0$，总余额为 $s$。任何修复至少插入 $L=-m$ 个左括号，否则最亏空前缀仍不合法。若插入左右括号数分别为 $\ell,r$，末尾平衡要求 $r=s+\ell$，故总插入数 $s+2\ell\ge s+2L$。算法恰插入 $L$ 个左括号，扫描后的余额为 $s+L$，再插入这么多个右括号，达到下界。时间 $O(T)$，除输出外只需常数空间。

\algproblem{ALG11}{二维地形的蓄水量}
矩形网格每格地面高度为非负实数 $h_v$，每格面积为1，相邻仅指上下左右。矩形外部可以自由排水，所有边界格都与外部相连。持续降雨至稳定后，求能够保留的总水量。给出算法并解释为什么只比较每行、每列最高点是不够的。
\solution 对每格 $v$ 定义最低逃逸水位
\[
\ell_v=\min_{\gamma:\,v\text{ 到边界}}\ \max_{u\in\gamma}h_u.
\]
路径包括起终点，故 $\ell_v\ge h_v$；某条路径的最高点决定沿它流出的门槛，而在所有逃逸路径中选最低门槛，得到稳定水位。总水量为 $\sum_v(\ell_v-h_v)$。

用最小堆计算这个“最小最大值”路径问题。所有边界格以距离 $h_v$ 入堆，内部距离初值为 $+\infty$。每次取距离最小的尚未确定格 $v$，对每个相邻格 $u$ 尝试松弛
\[
d_u\leftarrow\min\{d_u,\max(d_v,h_u)\}.
\]
忽略已过时的堆记录。由于沿路径扩展只会使最大高度不减，取出时未确定格不可能通过距离更大的路径得到更小值；这与最短路的贪心证明相同。更具体地，若有更低的逃逸路径，考察该路径上第一个未确定格，它应已由相邻的确定格获得不大的候选值，与当前最小堆顶矛盾。

网格边数为格数的常数倍，时间 $O(RC\log(RC+1))$，空间 $O(RC)$。只有一行或一列时每格都在边界，水量为零。实际逃逸路径可以转弯，因此只看直线方向的挡墙会高估蓄水。

\algproblem{ALG12}{全一正方形的最大面积与总数}
给定二元矩阵，同时求全为1的轴对齐正方形的最大面积，以及这种正方形的总个数。不同位置或边长均算不同正方形。要求 $O(RC)$ 时间、$O(C)$ 额外空间。
\solution 令 $d_{i,j}$ 为以 $(i,j)$ 作右下角的最大全一正方形边长。当前格为0时 $d_{i,j}=0$；为1时
\[
d_{i,j}=1+\min(d_{i-1,j},d_{i,j-1},d_{i-1,j-1}),
\]
越界项为0。若存在边长 $k$ 的正方形，则删掉末行、末列后，三个相邻位置都至少支持边长 $k-1$，给出必要性。反过来，三个边长 $k-1$ 的正方形与当前格的并集覆盖整个边长 $k$ 的方块，给出充分性。

最大面积是 $(\max d_{i,j})^2$。固定右下角时，可行边长恰为 $1,\ldots,d_{i,j}$，所以总数为
\[
\sum_{i,j}d_{i,j}.
\]
每个正方形有唯一右下角，因此没有重复计数。按行更新长度 $C$ 的数组时，额外保留被覆盖前的左上值；旧本列是上方，新前一列是左方。全零或空矩阵两项答案均为0。
